% ===================================================================
%  BAGAN No. 2 — Tubuh Manusia. --- Waktu Istirahat. Permainan.
%
%  Traduction, faite en 2026, en indonesien d'aujourd'hui, sur l'ido
%  de texto/io. Regles de la colonne : voir 10-bagan-01.tex. La cle
%  t02-tit-1 porte son « pos= » comme du cote ido : l'imprimeur pose
%  « LA KAPO. » avant l'alinea qui annonce les trois parties du corps,
%  et « pos= » ne deplace que la page de lecture.
%
%  L'ANATOMIE PASSE ENTIERE, SANS UN SEUL EMPRUNT, ET C'EST LA
%  TROISIEME COLONNE OU CELA ARRIVE. Kepala, tengkorak, otak, rambut,
%  wajah, dahi, pelipis, mata, alis, bulu mata, hidung, lubang hidung,
%  pipi, telinga, mulut, dagu, kumis, cambang, bibir, rahang, gigi,
%  lidah, leher, tenggorokan, tengkuk, dada, paru-paru, jantung,
%  tulang rusuk, punggung, bahu, perut, lambung, usus, lengan, tungkai,
%  otot, siku, pergelangan, tangan, jari, kuku, urat, darah, paha,
%  lutut, betis, daging, kulit, mata kaki, telapak, tumit :
%  cinquante-deux mots, tous malais de fond. L'ourdou avait le meme resultat
%  par un autre chemin — il a DEUX jeux complets, l'indien et l'arabe,
%  et choisit l'indien parce que le lecteur est un ecolier. Ici il n'y
%  a rien a choisir : la couche neerlandaise, si visible ailleurs dans
%  cette colonne, n'a pas touche au corps.
%
%  C'EST LA MEILLEURE MESURE DU RELEVE SUR CE POINT. Les trois
%  couches de l'indonesien se partagent le lexique par DOMAINE, non au
%  hasard : le corps et les gestes sont malais, l'ecole coloniale a
%  laisse le neerlandais (boven, grip, spons, sabak au tableau 1), et
%  l'abstrait vient de l'arabe et du sanskrit (guru, ilmu, sejarah).
%  Le corps est ce qu'on a avant toute ecole : il n'emprunte pas.
%
%  QUATRIEME REFUS AU MEME RENVOI, ET C'EST LE PLUS DUR DES QUATRE.
%  Le « baro-ludo » (38) — le jeu de barres — a fait trebucher toutes
%  les colonnes du releve : le coreen a refuse 진놀이, le tamoul
%  கபடி, l'ourdou کبڈی. L'indonesien a mieux encore : le gobak sodor,
%  aussi appele « hadang », est un jeu de lignes entre deux camps, et
%  il ressemble au jeu de barres BIEN PLUS que le kabaddi. La
%  tentation etait donc la plus forte des quatre. La colonne ecrit
%  « permainan barisan », qui decrit, et ne met pas un jeu javanais a
%  la place d'un jeu francais. On modernise la langue, jamais les
%  choses.
%
%  ET TROIS JEUX SE GARDENT, PARCE QUE CE SONT LES MEMES CHOSES.
%  gasing la toupie (18), kelereng les billes (21), egrang les
%  echasses (24), layang-layang le cerf-volant (33), petak umpet
%  cache-cache (25) : ce sont des mots indonesiens de fond pour des
%  objets et des regles identiques, pas des substitutions. La
%  difference d'avec le jeu de barres tient en une phrase : une toupie
%  est une toupie partout, un jeu de barres ne l'est pas.
%
%  bulu tangkis (29) DEMANDE UNE PHRASE A PART. L'ido decrit un volant
%  qui vole d'une raquette a l'autre — le jeu d'avant le badminton.
%  L'indonesien le nomme bulu tangkis, qui est aujourd'hui le sport
%  national du pays. Le mot est juste et la chose est la meme ; c'est
%  seulement le prestige qui a change de camp depuis 1926.
% ===================================================================

%%K t02-apar-1 apar
\VUsaut{4.0mm}
\VUtitre{15.0pt}{60}{BAGAN No. 2}
\VUsaut{2.0mm}
\VUtitre{11.4pt}{0}{Tubuh Manusia. --- Waktu Istirahat.}
\VUtitre{11.4pt}{0}{Permainan.}

%%K t02-tit-0 sub
\VUtitre{13.2pt}{0}{Tubuh manusia.}
\VUsaut{2.4mm}
\VUcentre{10.2pt}{0}{\textit{(Pelajaran sederhana ilmu alam.)}}
\VUsaut{3.0mm}

%%K t02-tit-1 sub pos=t02-01-1
\VUpk{10.2pt}{40}{Kepala.}
\VUsaut{3.0mm}

%%K t02-01-1 p
1. --- Bagian utama tubuh manusia adalah : \VUgras{kepala,}
\VUgras{badan} dan \VUgras{anggota badan.}

%%K t02-02-1 p
2. --- Kepala terdiri atas dua bagian : \VUgras{tengkorak,} yang berisi
\VUgras{otak} dan tertutup oleh \VUgras{rambut,} dan \VUgras{wajah.}

%%K t02-03-1 p
3. --- Pada gambar kita tidak melihat tengkorak maupun otak ; tetapi
kita melihat dengan sangat jelas bagian-bagian penting wajah.

%%K t02-04-1 p
4. --- Inilah pertama-tama \VUgras{dahi,} yang menandakan kecerdasan
yang kuat, pikiran yang selalu bekerja. \VUgras{Pelipisnya} sangat
lapang. \VUgras{Mata} itu hidup dan terbuka lebar. \VUgras{Alis} yang
tebal dan \VUgras{bulu mata} yang panjang melindunginya.

%%K t02-05-1 p
5. --- Inilah pula \VUgras{hidung} dan \VUgras{lubang hidung.} Hidung
berguna untuk mencium dan untuk bernapas. Di kedua sisi hidung ada
\VUgras{pipi,} dan lebih ke tepi ada \VUgras{telinga.} Bagian bawah
wajah terdiri atas \VUgras{mulut} dan \VUgras{dagu.}

%%K t02-06-1 p
6. --- Kita tidak melihatnya dengan baik, sebab \VUgras{kumis} dan
\VUgras{cambang} menutupinya dari pandangan kita. Tetapi kita tahu
bahwa mulut mempunyai dua \VUgras{bibir,} \VUgras{rahang atas} dan
\VUgras{rahang bawah} dengan \VUgras{giginya,} serta \VUgras{lidah.}
Dengan mulut kita berbicara dan makan. Dengan lidah kita mengecap
makanan.

%%K t02-07-1 p
7. --- Mata, telinga, hidung dan mulut, sama seperti jari-jari,
adalah alat dari kelima indra : \VUgras{indra penglihat,}
\VUgras{indra pendengar,} \VUgras{indra pencium,} \VUgras{indra
pengecap,} \VUgras{indra peraba.}

%%K t02-08-1 p
8. --- Jadi kepala adalah salah satu bagian tubuh manusia yang paling
menarik.

%%K t02-tit-2 sub
\VUsaut{4.4mm}
\VUpk{10.2pt}{40}{Badan.}
\VUsaut{3.0mm}

%%K t02-09-1 p
9. --- \VUgras{Leher} menghubungkan kepala dengan badan. Leher
meliputi \VUgras{tenggorokan} dan \VUgras{tengkuk.}

%%K t02-10-1 p
10. --- Badan juga berisi banyak organ yang penting. Di dalam
\VUgras{dada} ada \VUgras{paru-paru,} yang dengannya kita bernapas,
dan \VUgras{jantung,} yang merupakan pusat kehidupan. Bila jantung
berhenti berdenyut, makhluk hidup itu mati.

%%K t02-11-1 p
11. --- \VUgras{Tulang-tulang rusuk} menopang dada.

%%K t02-12-1 p
12. --- Bagian badan yang lain adalah \VUgras{pinggang,}
\VUgras{punggung,} \VUgras{bahu} dan \VUgras{perut.} Kita berbaring
pada pinggang atau pada punggung untuk tidur, dan kita memikul beban
di atas bahu kita.

%%K t02-13-1 p
13. --- Di dalam perut ada \VUgras{lambung} dan \VUgras{usus.}
Keduanya berguna untuk mencernakan makanan.

%%K t02-tit-3 sub
\VUsaut{4.4mm}
\VUpk{10.2pt}{40}{Anggota badan.}
\VUsaut{3.0mm}

%%K t02-14-1 p
14. --- Kita mempunyai empat \VUgras{anggota badan} : dua
\VUgras{lengan} dan dua \VUgras{tungkai.} Dengan lengan kita
mengambil, memegang, mengangkat dan memungut benda ; dengan tungkai
kita berdiri, berjalan dan berlari.

%%K t02-15-1 p
15. --- \VUgras{Bahu} menghubungkan lengan dengan tubuh. Sebuah
\VUgras{otot} yang sangat kuat, bisep, menggerakkan lengan, dan
\VUgras{siku} menghubungkannya dengan \VUgras{lengan bawah.}
\VUgras{Pergelangan} membentuk sendi \VUgras{tangan,} yang berakhir
pada \VUgras{jari-jari} dan \VUgras{kuku.} Pada tangan kita
membedakan \VUgras{urat} (dan nadi) yang di dalamnya mengalir
\VUgras{darah.}

%%K t02-16-1 p
16. --- Bagian tungkai adalah : \VUgras{paha,} \VUgras{lutut} dan
\VUgras{lipat lutut,} \VUgras{betis,} yang \VUgras{dagingnya}
tertutup oleh \VUgras{kulit,} \VUgras{mata kaki} dan \VUgras{kaki.}
\VUgras{Tulang} tungkai sangat besar dan sangat kuat. Di ujung kaki
ada \VUgras{jari-jari kaki.} Bagian selebihnya adalah \VUgras{telapak}
dan \VUgras{tumit.}

%%K t02-17-1 p
17. --- Demikianlah gambaran tubuh manusia yang hampir lengkap.

%%K t02-17-2 p
\textit{Catatan.} --- \textit{Dengan bantuan ungkapan « saya sakit
pada... (kepala dan sebagainya) », guru akan dapat memakai kembali
seluruh kosakata ini dari segi baik atau buruknya} \VUgras{keadaan
tubuh.}

%%K t02-tit-4 sub
\VUsaut{5.0mm}
\VUpk{10.2pt}{40}{Senam.}
\VUsaut{2.4mm}
\VUcentre{10.2pt}{0}{\textit{(Cerita salah seorang murid.)}}
\VUsaut{3.0mm}

%%K t02-18-1 p
18. --- Supaya lentur, tangkas dan sehat, kita harus melatih tungkai
dan lengan kita. Karena itu kita bermain dan berlatih \VUgras{senam.}

%%K t02-19-1 p
19. --- Selama hari sekolah, kedua latihan itu berlangsung di halaman
sekolah menengah (lihat bagan No. 1). Tetapi pada hari Kamis dan hari
Minggu kami pergi ke lapangan rumput yang luas, tempat kami punya
ruang untuk berlari dan bermain.

%%K t02-20-1 p
20. --- Guru-guru dan orang tua kami kadang-kadang menemani kami ke
sana ; tetapi yang memimpin latihan tentulah \VUgras{guru
senam}\textsuperscript{(12)}.

%%K t02-21-1 p
21. --- Di lapangan rumput, di sebelah kiri pintu masuk, ada
\VUgras{alat-alat senam}\textsuperscript{(1)} : kami memanjat
\VUgras{tangga}\textsuperscript{(2)}, kami berjungkir balik pada
\VUgras{gelang-gelang}\textsuperscript{(3)} dan pada
\VUgras{trapeze}\textsuperscript{(4)}, kami menarik diri dengan tangan
sampai ke puncak \VUgras{tangga tali}\textsuperscript{(5)} atau
\VUgras{tali bersimpul}\textsuperscript{(6)}.

%%K t02-22-1 p
22. --- Sementara itu, kawan-kawan kami melompati \VUgras{kuda-kuda
kayu}\textsuperscript{(7)} atau \VUgras{palang
sejajar}\textsuperscript{(11)}. Yang lain \VUgras{bertinju}\textsuperscript{(8)}
atau \VUgras{bermain anggar}\textsuperscript{(9)} dengan topeng dan
\VUgras{floret.} Yang lain lagi \VUgras{mengangkat
halter}\textsuperscript{(10)}.

%%K t02-tit-5 sub
\VUsaut{4.6mm}
\VUpk{10.2pt}{40}{Permainan.}
\VUsaut{3.0mm}

%%K t02-23-1 p
23. --- Di sekitar para pesenam, anak laki-laki dan anak perempuan
bermain bermacam-macam \VUgras{permainan.} Anak-anak perempuan
bersenang-senang dengan \VUgras{simpainya}\textsuperscript{(14)} ;
mereka melompati \VUgras{tali}\textsuperscript{(13)} atau mengajak
saudara dan sepupu laki-laki mereka bermain
\VUgras{kroket}\textsuperscript{(15)}. Mereka senang sekali mendorong
pergi \VUgras{bola} lawan dengan \VUgras{palu kayu} mereka, tepat pada
saat lawan itu mengira akan lewat dengan mudah di bawah gawang kecil!

%%K t02-24-1 p
24. --- Anak-anak laki-laki lebih suka permainan yang lebih menguras
otot. Mereka senang bermain \VUgras{gada-gada}\textsuperscript{(16)},
yang mereka robohkan dengan cekatan memakai
\VUgras{bola}\textsuperscript{(17)}. Mereka juga tidak menolak bermain
\VUgras{gasing}\textsuperscript{(18)} dan
\VUgras{bilbuket}\textsuperscript{(19)} atau bahkan \VUgras{permainan
katak}\textsuperscript{(20)}. Tetapi permainan kesukaan mereka
tentulah \VUgras{kelereng}\textsuperscript{(21)} dan \VUgras{bola
tembok}\textsuperscript{(22)}, sebab dinding \VUgras{rumah
kaca}\textsuperscript{(26)} sangat cocok untuk memantulkan bola ringan
itu.

%%K t02-25-1 p
25. --- Si kecil Ioannes, yang berjalan dengan
\VUgras{egrangnya}\textsuperscript{(24)}, sangat cekatan dan pandai
sekali menjaga keseimbangan. Di dekatnya, beberapa kawan bermain
\VUgras{petak umpet}\textsuperscript{(25)} di rumah kaca itu dan di
balik \VUgras{pagar tanaman.} Yang lain lebih suka \VUgras{tebak
pukulan}\textsuperscript{(28)} atau \VUgras{bulu
tangkis}\textsuperscript{(29)}, yang koknya melayang dari satu
\VUgras{raket} ke raket yang lain.

%%K t02-noto-1 noto
\VUnotes{143.2mm}{%
(*) Permainan anak-anak sering mempunyai nama yang murni kebangsaan.
Kami berusaha menamainya dalam Ido sebaik mungkin — kadang dengan
mengambil ilham dari perbuatan yang menandainya, kadang dengan memilih
nama kebangsaan negeri ini atau itu, bila nama itu tidak terlalu tidak
tepat. Misalnya : \VUgras{permainan tong} (di Jerman dan Prancis)
adalah \VUgras{permainan katak} \textsuperscript{(20)} (dalam bahasa
Inggris), yang di dalamnya para pemain berusaha melemparkan cakram
bundar mereka melalui mulut katak logam yang berdiri menganga di atas
peti (tong) berlubang.
\VUgras{Lompat bahu}\textsuperscript{(30)} berbeda dari \VUgras{lompat
punggung} hanya dalam hal ini : untuk melompati kepala, para pemain
menumpukan tangan pada bahu pasangannya, sedangkan dalam «
\VUgras{lompat punggung} » mereka menumpukan tangan hanya pada
punggungnya. --- Atau juga : sebagian peserta membungkuk sementara
yang lain akan melompati punggung mereka dengan menumpukan tangan pada
bahu ; yang pertama harus menahan benturan tanpa melengkung, yang
kedua harus melompat tanpa kehilangan keseimbangan (lihat bagan itu
sendiri).%
}

%%K t02-26-1 p
26. --- Di tengah lapangan rumput ada dua pohon. Di kaki salah
satunya, tujuh atau delapan anak laki-laki mengadakan permainan
\VUgras{lompat bahu}\textsuperscript{(30)} (*). Tidak di semua negeri
permainan itu dikenal ; tetapi semua orang tahu apa itu \VUgras{lompat
punggung,} yang kali ini tidak mereka mainkan.

%%K t02-tit-6 sub
\VUsaut{5.0mm}
\VUpk{10.2pt}{40}{Permainan di udara terbuka.}
\VUsaut{3.4mm}

%%K t02-27-1 p
27. \VUgras{Permainan mata tertutup}\textsuperscript{(31)} juga sangat
menyenangkan ; anak perempuan dan anak laki-laki bergantian memasang
\VUgras{penutup mata} pada mata mereka dan menangkap yang kurang
cekatan, yang membiarkan dirinya tertangkap.

%%K t02-28-1 p
28. --- Di tengah lapangan rumput, inilah permainan yang sangat
Prancis tetapi agak keras : itulah \VUgras{permainan
beruang}\textsuperscript{(32)}, jauh lebih ribut daripada permainan
\VUgras{layang-layang}\textsuperscript{(33)} yang begitu tenang, atau
bahkan daripada permainan Inggris
\VUgras{tenis}\textsuperscript{(34)}.

%%K t02-29-1 p
29. --- Olahraga yang terakhir itu adalah kegembiraan murid-murid
besar. Guru-guru kami sendiri dengan senang hati memainkannya. Ia
menuntut kecakapan besar dengan ketangkasan, dan ia sangat
menyenangkan bagi saya. Permainan \VUgras{ayunan}\textsuperscript{(35)}
saya tinggalkan untuk anak-anak perempuan.

%%K t02-30-1 p
30. --- Saya tidak menyukai permainan Inggris yang lain, yang barangkali
bisa menjadi berbahaya.

%%K t02-30-1b p
Yang saya maksud adalah \VUgras{sepak bola}\textsuperscript{(36)}, yang
menggembirakan kakak laki-laki saya. Saya lebih suka
\VUgras{permainan barisan}\textsuperscript{(38)} kami yang lama, yang
sama sehatnya dan sungguh kurang kasar.

%%K t02-tit-7 sub
\VUsaut{4.6mm}
\VUpk{10.2pt}{40}{Bersepeda dan bermain balon.}
\VUsaut{3.0mm}

%%K t02-31-1 p
31. --- Dengan senang hati pula saya mengelilingi lapangan rumput
dengan \VUgras{sepeda}\textsuperscript{(37)}.

%%K t02-32-1 p
32. --- Tahun depan saya akan belajar \VUgras{menunggang
kuda}\textsuperscript{(39)}. Betapa indahnya kuda yang melangkah atau
berlari kecil dengan anggun, dan yang berderap melompati rintangan!

%%K t02-33-1 p
33. --- Pada musim panas kami belajar \VUgras{berenang}\textsuperscript{(40)}.
Dengan senang hati kami menyelam dan mandi di air sungai yang jernih
dan segar. Kadang-kadang, akhirnya, kami menerbangkan
\VUgras{balon}\textsuperscript{(41)} kertas yang naik dengan megah ke
udara, begitu ia diisi dengan udara panas.

%%K t02-34-1 p
34. --- Masih ada banyak permainan lain, tetapi saya tidak akan
selesai menyebutkannya, dan dari ringkasan ini orang mengerti bahwa
pada hari Kamis kami tidak terlalu bosan di lapangan rumput kami yang
indah.

%%K t02-34-2 p
N.-B. --- \textit{Guru akan dengan mudah melengkapi bab ini, dengan
menggambarkan lebih terperinci masing-masing permainan yang paling
penting.}

\VUsaut{6.0mm}
\VUfilet{22mm}
